% ==============================================================================
% Chapitre 4 : Infrastructure et Déploiement
% Docker, CI/CD, scripts, monitoring, sécurité
% ==============================================================================

\chapter{Infrastructure et Déploiement}

Ce chapitre décrit l'infrastructure technique qui supporte l'ensemble du
système LockBits. Il couvre la conteneurisation avec Docker, le pipeline
d'intégration et de déploiement continus avec GitHub Actions, les scripts
d'installation et de maintenance, la stack de monitoring, ainsi que les
mesures de sécurité mises en œuvre.

\section{Docker et Conteneurisation}

\subsection{Stratégie de conteneurisation}

La conteneurisation est au cœur de l'architecture LockBits. Chaque composant du
système possède son propre \texttt{Dockerfile} et est empaqueté dans une image
Docker autonome. Cette approche garantit la reproductibilité des déploiements
et la portabilité entre les environnements de développement, de pré-production
et de production.

Les images sont construites avec \textbf{Docker Buildx} pour supporter deux
architectures matérielles~: \texttt{linux/amd64} (x86\_64) et
\texttt{linux/arm64} (AArch64). Cette compatibilité multi-architecture permet
de déployer la stack aussi bien sur des serveurs x86 traditionnels que sur des
instances ARM plus économes en énergie (AWS Graviton, Raspberry Pi pour les
tests). Une fois construites, les images sont publiées sur le \textbf{GitHub
Container Registry (GHCR)} avec plusieurs Tags~:

\begin{itemize}
  \item \texttt{latest}~: dernier build validé de la branche \texttt{main}.
  \item \texttt{develop}~: dernier build de la branche \texttt{develop}.
  \item \texttt{vX.Y.Z}~: version sémantique pour les releases officielles.
\end{itemize}

\begin{table}[H]
\centering
\caption{Images Docker du projet LockBits}
\label{tab:images-docker}
\begin{tabularx}{\textwidth}{l l l X}
\toprule
\textbf{Image GHCR} & \textbf{Composant} & \textbf{Technologie} & \textbf{Description} \\
\midrule
\texttt{ghcr.io/lockbits-esgi/edr} & Serveur EDR & Python / FastAPI & API REST de détection et réponse \\
\texttt{ghcr.io/lockbits-esgi/edr-agent} & Agent EDR & Python / PyInstaller & Agent de collecte déployable \\
\texttt{ghcr.io/lockbits-esgi/site\_lockbits} & Portail client & PHP 8+ / Apache & Site web de gestion clients \\
\bottomrule
\end{tabularx}
\end{table}

\subsection{Dockerfiles multi-stage}

Les \texttt{Dockerfile} utilisent une construction en deux étapes (multi-stage
build) pour réduire la taille des images finales. La première étape (\emph{builder})
installe les dépendances de compilation, tandis que la seconde étape (\emph{final})
ne contient que les dépendances d'exécution et le code applicatif.

Extrait du Dockerfile du serveur EDR~:

\begin{lstlisting}[language=dockerbash, caption={Dockerfile multi-stage du serveur EDR}, label={lst:dockerfile-edr}]
# Stage 1: Builder
FROM python:3.14-slim AS builder
WORKDIR /build
RUN apt-get update && apt-get install -y --no-install-recommends gcc \\
    && rm -rf /var/lib/apt/lists/*
COPY requirements.txt .
RUN pip install --no-cache-dir --user -r requirements.txt

# Stage 2: Final
FROM python:3.14-slim
WORKDIR /app
RUN apt-get update && apt-get upgrade -y && apt-get install -y \\
    --no-install-recommends curl \\
    && rm -rf /var/lib/apt/lists/*
COPY --from=builder /root/.local /root/.local
ENV PATH=/root/.local/bin:\$PATH
COPY . .
EXPOSE 8000
HEALTHCHECK --interval=30s --timeout=5s --start-period=10s --retries=3 \\
    CMD curl -f http://localhost:8000/health || exit 1
CMD ["python", "-m", "uvicorn", "server.main:app", "--host", "0.0.0.0", "--port", "8000"]
\end{lstlisting}

Cette approche présente plusieurs avantages. La taille de l'image finale est
réduite, ce qui accélère les téléchargements et diminue la surface d'attaque.
Les dépendances sont figées dans l'image, garantissant un environnement
d'exécution identique en développement et en production. Enfin, le healthcheck
intégré permet à Docker de surveiller l'état du service et de redémarrer
automatiquement le conteneur en cas de défaillance.

\subsection{Docker Compose}

L'orchestration des services est gérée par Docker Compose. Trois fichiers de
composition couvrent les différents environnements~:

\begin{itemize}
  \item \textbf{\texttt{docker-compose.yml}}~: environnement de développement
    local. Les images sont construites depuis les sources (directive
    \texttt{build}) avec rechargement à chaud.
  \item \textbf{\texttt{docker-compose.prod.yml}}~: environnement de
    production. Les images sont tirées depuis GHCR avec le tag \texttt{:latest}
    et la politique de téléchargement \texttt{pull\_policy: always}.
  \item \textbf{\texttt{docker-compose.preprod.yml}}~: environnement de
    pré-production. Les images utilisent le tag \texttt{:develop}, avec des
    noms de conteneurs suffixés par \texttt{\_preprod} pour éviter les conflits
    avec la stack de production.
\end{itemize}

Chaque fichier de composition définit un \texttt{project name} explicite
(\texttt{lockbits-prod} ou \texttt{lockbits-preprod}) pour garantir l'isolation
des réseaux Docker entre les environnements. Sans cette précaution, les
réseaux partagés (notamment le réseau \texttt{monitoring}) provoqueraient un
round-robin DNS entre les instances de production et de pré-production.

Extrait du fichier \texttt{docker-compose.prod.yml}~:

\begin{lstlisting}[language=dockercompose, caption={Extrait du docker-compose.prod.yml -- service EDR}, label={lst:docker-compose-prod}]
services:
  edr-server:
    image: ghcr.io/lockbits-esgi/edr:latest
    container_name: lockbits_edr
    pull_policy: always
    depends_on:
      edr-db:
        condition: service_healthy
    ports:
      - "${EDR_PORT:-8001}:8000"
    environment:
      SERVER_HOST: "0.0.0.0"
      SERVER_PORT: "8000"
      DATABASE_URL: ${EDR_DATABASE_URL:-postgresql://${EDR_DB_USER:-edr}:${EDR_DB_PASS}@edr-db/${EDR_DB_NAME:-edr}}
      AUTH_TOKEN: ${EDR_AUTH_TOKEN}
      VT_ENABLED: ${VT_ENABLED:-false}
    volumes:
      - edr_data:/app/data
    networks:
      - frontend
      - backend
      - monitoring
    restart: unless-stopped
    deploy:
      resources:
        limits:
          cpus: "1"
          memory: 512M
\end{lstlisting}

Le service \texttt{edr-server} utilise une base PostgreSQL hébergée dans un
conteneur dédié (\texttt{edr-db}) avec des limites de ressources explicites.
Les variables d'environnement critiques (mot de passe de base, token d'auth)
sont passées via le fichier \texttt{.env} et n'apparaissent jamais en dur dans
la configuration.

\subsection{Réseau Docker}

L'infrastructure réseau est segmentée en trois réseaux bridge distincts,
chacun avec un niveau d'exposition différent~:

\begin{itemize}
  \item \textbf{Réseau \texttt{frontend}}~: exposé sur l'hôte, il connecté le
    serveur EDR et le site web. Seuls ces deux services sont accessibles
    depuis l'extérieur via les ports publiés.
  \item \textbf{Réseau \texttt{backend}}~: réseau interne qui connecté le
    serveur EDR à sa base de données PostgreSQL. Aucun port n'est exposé sur
    l'hôte pour ce réseau.
  \item \textbf{Réseau \texttt{monitoring}}~: réseau interne dédié aux
    services de supervision (Prometheus, Grafana, Loki, Alloy). Les services
    applicatifs y sont connectés pour exposer leurs métriques, mais le réseau
    lui-même n'est pas directement accessible depuis l'extérieur.
\end{itemize}

\begin{table}[H]
\centering
\caption{Ports exposés par les services LockBits}
\label{tab:ports-exposes}
\begin{tabularx}{\textwidth}{l l l X}
\toprule
\textbf{Service} & \textbf{Port hôte} & \textbf{Port conteneur} & \textbf{Usage} \\
\midrule
Serveur EDR & 8001 & 8000 & API REST EDR \\
Site web & 8080 & 80 & Portail client HTTP \\
Prometheus & 9092 / 9090 & 9090 & Métriques (prod / preprod) \\
Grafana & 3002 / 3001 & 3000 & Tableaux de bord (prod / preprod) \\
\bottomrule
\end{tabularx}
\end{table}

\section{CI/CD avec GitHub Actions}

\subsection{Pipeline CI}

Le pipeline d'intégration continue est défini dans le fichier
\texttt{.github/workflows/build-and-push.yml}, un workflow entièrement
\textbf{homemade} (fait maison). Il se déclenche sur chaque push sur les
branches \texttt{main} et \texttt{develop}, ainsi que sur les pull requests
vers ces branches. Le pipeline se déroule en quatre étapes principales~:

\begin{enumerate}
  \item \textbf{Lint}~: vérification du code Python du serveur EDR avec Ruff
    (lint et format), exécutée en parallèle des tests.
  \item \textbf{Tests}~: tests unitaires avec pytest pour le serveur EDR,
    PHPUnit pour le site web, et tests d'intégration (smoke tests) avec des
    conteneurs Docker lancés dans le runner CI.
  \item \textbf{Build et push}~: construction des images Docker multi-architecture
    avec Docker Buildx, puis publication sur GHCR. Cette étape ne s'exécute
    que pour les pushes (pas pour les pull requests).
  \item \textbf{Scan de sécurité}~: analyse des images publiées avec Trivy
    pour détecter les vulnérabilités CRITICAL et HIGH.
\end{enumerate}

\begin{lstlisting}[language=yaml, caption={Extrait du pipeline CI -- tests unitaires EDR}, label={lst:ci-edr-tests}]
jobs:
  lint-edr:
    name: "EDR - Ruff lint + format"
    runs-on: ubuntu-latest
    steps:
      - uses: actions/checkout@v6
        with:
          submodules: recursive
      - uses: actions/setup-python@v6
        with:
          python-version: "3.12"
      - name: Install Ruff
        run: pip install ruff
      - name: Ruff check
        working-directory: edr
        run: ruff check --output-format=github
      - name: Ruff format check
        working-directory: edr
        run: ruff format --check

  test-edr-unit:
    name: "EDR - Unit tests"
    runs-on: ubuntu-latest
    steps:
      - uses: actions/checkout@v6
        with:
          submodules: recursive
      - uses: actions/setup-python@v6
        with:
          python-version: "3.12"
          cache: pip
          cache-dependency-path: edr/requirements.txt
      - name: Install dependencies
        run: pip install -r edr/requirements.txt
      - name: Run pytest
        working-directory: edr
        run: pytest tests/ -v --tb=short
        env:
          DATABASE_URL: "sqlite:///:memory:"
          VT_ENABLED: "false"
\end{lstlisting}

La construction des images est réalisée à l'aide de l'action
\texttt{docker/build-push-action} avec la prise en charge de la cache
GitHub Actions (\texttt{cache-from}/\texttt{cache-to}) pour accélérer les
builds successifs.

\subsection{Mise à jour automatique des submodules}

Le dépôt principal orchestre les composants via des sous-modules Git. Un
workflow dédié (\texttt{auto-update-submodules.yml}) vérifie périodiquement
si les sous-modules ont de nouveaux commits. Si c'est le cas, il met à jour
la référence, crée un commit et le pousse sur la branche \texttt{develop}, ce
qui déclenche automatiquement le pipeline CI et le rebuild des images.

\begin{lstlisting}[language=yaml, caption={Workflow de mise à jour automatique des submodules}, label={lst:auto-update-submodules}]
name: Auto-update submodules

on:
  schedule:
    - cron: "0 6 * * *"
  workflow_dispatch:

jobs:
  update:
    name: Update submodules
    runs-on: ubuntu-latest
    permissions:
      contents: write
    steps:
      - uses: actions/checkout@v6
        with:
          ref: develop
          submodules: recursive
          token: ${{ secrets.GITHUB_TOKEN }}

      - name: Update submodules to latest remote
        run: |
          git fetch --deepen=100 origin 2>/dev/null || true
          git submodule update --remote

      - name: Check if anything changed
        id: diff
        run: |
          if git diff --quiet --exit-code; then
            echo "changed=false" >> $GITHUB_OUTPUT
          else
            echo "changed=true" >> $GITHUB_OUTPUT
          fi

      - name: Commit and push
        if: steps.diff.outputs.changed == 'true'
        run: |
          git config user.name "LockBits Bot"
          git config user.email "bot@lockbits-esgi"
          git add -u
          git commit -m "chore: auto-update submodules"
          git push origin develop
\end{lstlisting}

Ce mécanisme entièrement automatisé permet de maintenir les sous-modules à jour
sans intervention humaine. La synchronisation a lieu une fois par jour à 6h
UTC, ce qui correspond à un rythme raisonnable compte tenu de la fréquence des
mises à jour des composants.

\subsection{Déploiement automatisé (Dokploy)}

Après la construction et la publication des images, le pipeline déclenche le
déploiement via un webhook \textbf{Dokploy}. Deux webhooks sont configurés~:

\begin{itemize}
  \item Un webhook pour la production (branche \texttt{main}) qui redéploie la
    stack avec les images \texttt{:latest}.
  \item Un webhook pour la pré-production (branche \texttt{develop}) qui
    redéploie la stack avec les images \texttt{:develop}.
\end{itemize}

Le déclenchement est réalisé par une requête HTTP POST vers l'URL du webhook
configurée dans les \texttt{secrets} GitHub. Les tokens Dokploy sont stockés
dans les \textbf{secrets GitHub} (\texttt{DOKPLOY\_PROD\_HOOK} et
\texttt{DOKPLOY\_PREPROD\_HOOK}) et ne sont jamais exposés dans le code source.

\subsection{Domaines et URLs publiques}

L'ensemble des services est déployé publiquement sur le domaine
\texttt{lockbits.pro}. Chaque composant possède son propre sous-domaine~:

\begin{itemize}
  \item \url{https://lockbits.pro} -- Site client (portail web)
  \item \url{https://edr.lockbits.pro} -- Serveur EDR (API REST)
  \item \url{https://preprod-edr.lockbits.pro} -- Serveur EDR (préproduction, branche \texttt{dev})
  \item \url{https://chat.lockbits.pro} -- Chatbot client
\end{itemize}

Tous les dépôts sont publics sur \url{https://github.com/Lockbits-ESGI}.

\subsection{Cloudflare Zero Trust}

Tous les services déployés sont protégés et exposés via \textbf{Cloudflare Zero
Trust}. L'infrastructure s'appuie sur un tunnel \texttt{cloudflared} exécuté
dans un conteneur LXC dédié, qui agit comme point d'entrée unique pour le
trafic externe.

\subsubsection{Architecture du tunnel}

Le tunnel \texttt{cloudflared} établit une connexion sortante depuis
l'infrastructure vers le réseau Cloudflare. Aucun port n'est ouvert
directement sur l'hébergement~; tout le trafic transite par le tunnel :

\begin{itemize}
  \item Le LXC \texttt{cloudflared} héberge le démon \texttt{cloudflared}
    configuré avec un tunnel permanent ;
  \item Cloudflare agit comme proxy inverse et terminatrice SSL/TLS ;
  \item Les requêtes sont routées vers les services internes selon les
    règles définies dans le tableau de bord Cloudflare ;
  \item Seul le trafic authentifié et filtré par Cloudflare atteint les
    services.
\end{itemize}

\subsubsection{Avantages}

\begin{itemize}
  \item \textbf{HTTPS automatique}~: Cloudflare fournit et renouvelle
    automatiquement les certificats SSL/TLS pour tous les sous-domaines ;
  \item \textbf{Protection DDoS}~: le réseau Cloudflare absorbe et filtre
    les attaques avant qu'elles n'atteignent l'infrastructure ;
  \item \textbf{AuthZero}~: possibilité d'activer l'authentification
    Zero Trust (email, OTP, SSO) pour restreindre l'accès à certains
    services (ex~: Grafana, monitoring) ;
  \item \textbf{Cache et performance}~: Cloudflare met en cache les
    ressources statiques et accélère la livraison du contenu ;
  \item \textbf{IP source masquée}~: l'adresse IP réelle de
    l'infrastructure n'est jamais exposée publiquement.
\end{itemize}

\subsubsection{Domaines couverts}

Les certificats SSL/TLS couvrent l'ensemble des sous-domaines
\texttt{lockbits.pro}~:
\texttt{lockbits.pro}, \texttt{edr.lockbits.pro},
\texttt{preprod-edr.lockbits.pro}, \texttt{chat.lockbits.pro}.

\section{Scripts d'installation et maintenance}

\subsection{install.sh}

Le script \texttt{install.sh} est le point d'entrée principal pour le
déploiement de la stack LockBits. Il s'exécute en une seule commande sans
nécessiter de clonage du dépôt~:

\begin{lstlisting}[language=dockerbash, caption={Commande d'installation one-liner}, label={lst:install-oneliner}]
curl -fsSL https://raw.githubusercontent.com/Lockbits-ESGI/main/main/install.sh | bash
\end{lstlisting}

Le script supporte plusieurs modes d'exécution~:

\begin{itemize}
  \item \textbf{Mode par défaut}~: installation complète de la stack. Le script
    télécharge les fichiers de configuration (\texttt{docker-compose.prod.yml}
    et \texttt{.env.example}), vérifie que Docker est installé, crée le fichier
    \texttt{.env}, tente de se connecter à GHCR (si un \texttt{GITHUB\_TOKEN}
    est fourni), puis lance la stack avec \texttt{docker compose up -d}.
  \item \texttt{--update}~: mise à jour rapide. Le script détecte les nouvelles
    variables d'environnement, redémarre la stack avec les dernières images et
    valide le bon fonctionnement des services.
  \item \texttt{--cron}~: installation d'une tâche cron pour la mise à jour
    automatique. Par défaut toutes les 5 minutes (configurable avec
    \texttt{--interval=N}).
  \item \texttt{--env=preprod}~: déploiement en pré-production avec le fichier
    \texttt{docker-compose.preprod.yml}.
  \item \texttt{--update-sources}~: mise à jour des références des sous-modules
    Git (depuis un clone du dépôt), suivie d'un commit et d'un push pour
    déclencher la CI.
\end{itemize}

Le script intègre également une validation de sécurité qui vérifie que les mots
de passe par défaut ont été remplacés avant le déploiement en production, et
refuse de continuer si ce n'est pas le cas.

\subsection{install-agent.sh}

Le script \texttt{install-agent.sh} est conçu pour le déploiement de l'agent
EDR sur les machines des clients. Il peut fonctionner selon deux modes~:

\begin{itemize}
  \item \textbf{Mode binaire}~: télécharge le binaire standalone (compilé avec
    PyInstaller) depuis les \textbf{GitHub Releases} et l'installe dans
    \texttt{/usr/local/bin/}. Un service systemd est créé pour assurer le
    démarrage automatique au boot.
  \item \textbf{Mode Docker}~: télécharge l'image GHCR de l'agent et configure
    un conteneur Docker avec les droits nécessaires (mode privilégié, accès au
    PID de l'hôte, montage du système de fichiers en lecture seule).
\end{itemize}

\begin{lstlisting}[language=dockerbash, caption={Installation de l'agent EDR (mode binaire)}, label={lst:install-agent}]
curl -fsSL https://raw.githubusercontent.com/Lockbits-ESGI/main/main/install-agent.sh | \\
  SERVER_URL=https://edr.lockbits.pro AUTH_TOKEN=token bash
\end{lstlisting}

L'agent supporte les architectures \texttt{amd64} et \texttt{arm64}. La
détection se fait automatiquement au moment de l'installation. Un fichier de
configuration YAML est généré dans \texttt{/etc/lockbits-agent/config.yaml}
avec les paramètres de connexion, la configuration du File Integrity Monitoring
(FIM) et le niveau de journalisation.

\subsection{backup.sh / restore.sh}

Les scripts \texttt{backup.sh} et \texttt{restore.sh} assurent la sauvegarde et
la restauration des données de l'infrastructure.

\texttt{backup.sh} crée des sauvegardes horodatées contenant~:

\begin{itemize}
  \item La base de données MySQL (dump via \texttt{mysqldump} avec les options
    \texttt{--single-transaction}, \texttt{--routines}, \texttt{--triggers}).
  \item La base de données PostgreSQL du serveur EDR (dump via \texttt{pg\_dump}
    avec \texttt{--clean} et \texttt{--if-exists}).
  \item En l'absence de PostgreSQL, les données du volume EDR sont sauvegardées
    via \texttt{docker cp}.
  \item Les fichiers de configuration (\texttt{.env}, \texttt{docker-compose*.yml}).
\end{itemize}

La rotation des sauvegardes conserve les 7 derniers jours. Le script
\texttt{restore.sh} liste les sauvegardes disponibles, valide leur contenu,
puis restaure chaque composant de manière interactive après confirmation de
l'utilisateur.

\section{Monitoring et Observabilité}

\subsection{Stack technique}

La supervision de l'infrastructure repose sur la stack \textbf{PLA}
(Prometheus, Loki, Alloy, Grafana)~:

\begin{itemize}
  \item \textbf{Prometheus}~: collecte les métriques des services exposant un
    endpoint \texttt{/metrics}. Chaque service applicatif (EDR, site web) est
    configuré comme cible de scraping.
  \item \textbf{Grafana}~: visualisation des métriques via des dashboards
    pré-configurés. L'authentification est requise en production, avec des
    identifiants administrateur définis dans les variables d'environnement.
  \item \textbf{Loki}~: agrégation centralisée des logs applicatifs. Il stocke
    les logs sous forme compressée et les rend interrogeables via Grafana.
  \item \textbf{Alloy}~: collecteur de logs Docker qui lit les fichiers JSON
    des conteneurs sur le système hôte et les achemine vers Loki.
\end{itemize}

\begin{lstlisting}[language=yaml, caption={Configuration Prometheus pour le scraping des services}, label={lst:prometheus-config}]
global:
  scrape_interval: 15s
  evaluation_interval: 15s

scrape_configs:
  - job_name: "lockbits_edr"
    metrics_path: "/metrics"
    static_configs:
      - targets: ["lockbits_edr:8000"]

  - job_name: "lockbits_site"
    metrics_path: "/metrics"
    static_configs:
      - targets: ["lockbits_web:80"]
\end{lstlisting}

Les cibles de scraping utilisent les noms de conteneurs (\texttt{lockbits\_edr},
\texttt{lockbits\_web}) plutôt que les noms de services Docker Compose
(\texttt{edr-server}, \texttt{site-web}). Cette précaution évite le round-robin
DNS qui se produirait si les deux stacks (production et pré-production)
partageaient le même réseau Dokploy.

\subsection{Métriques collectées}

Le système expose trois catégories de métriques~:

\begin{itemize}
  \item \textbf{Métriques applicatives EDR}~: nombre de requêtes API, temps de
    réponse, événements de sécurité traités, nombre d'agents actifs, taille de
    la file d'attente de rapports.
  \item \textbf{Métriques système}~: utilisation CPU, mémoire, disque et réseau
    des conteneurs, collectées via cAdvisor.
  \item \textbf{Métriques Docker}~: état des conteneurs (running, stopped,
    unhealthy), nombre de redémarrages, âge des conteneurs.
\end{itemize}

Les dashboards Grafana sont provisionnés automatiquement via le répertoire
\texttt{docker/grafana/provisioning/} qui contient les définitions des sources
de données et des tableaux de bord au format JSON.

\section{Sécurité de l'infrastructure}

\subsection{Scan Trivy}

Chaque image Docker est analysée par \textbf{Trivy} après sa publication sur
GHCR. Le scan recherche les vulnérabilités de niveau CRITICAL et HIGH, tant au
niveau du système d'exploitation que des bibliothèques logicielles. Le scan est
configuré en mode rapport (\texttt{exit-code: "0"})~: il ne bloque pas le
déploiement mais produit un rapport visible dans l'interface GitHub Actions.
Cette approche permet aux développeurs de prendre connaissance des
vulnérabilités sans interrompre le cycle de déploiement.

\subsection{Gestion des secrets}

La gestion des secrets suit les pratiques de sécurité suivantes~:

\begin{itemize}
  \item \textbf{Fichier \texttt{.env}}~: toutes les variables sensibles (mots
    de passe, tokens API, clés d'authentification) sont stockées dans un
    fichier \texttt{.env} qui n'est jamais commité (protégé par
    \texttt{.gitignore}). Le fichier \texttt{.env.example} fournit les noms des
    variables avec des valeurs par défaut.
  \item \textbf{Secrets GitHub}~: les tokens nécessaires à la CI/CD (GHCR,
    webhooks Dokploy) sont stockés dans les secrets du dépôt GitHub et injectés
    dans les workflows via la syntaxe \texttt{\${{ secrets.\textless{}nom\textgreater{}}}}.
  \item \textbf{Fichiers exclus}~: le \texttt{.gitignore} protège les fichiers
    sensibles~: \texttt{.env}, \texttt{miniedr.db}, \texttt{dist/},
    \texttt{build/}, \texttt{venv/}.
\end{itemize}

\subsection{CORS et rate limiting}

L'API REST du serveur EDR est protégée par deux mécanismes~:

\begin{itemize}
  \item \textbf{CORS}~: la configuration \texttt{CORS\_ORIGINS} permet de
    restreindre les origines autorisées à accéder à l'API. En développement,
    la valeur par défaut \texttt{*} autorise toutes les origines. En
    production, elle doit être remplacée par l'URL du portail client.
  \item \textbf{Rate limiting}~: les endpoints de l'API EDR sont protégés par
    une limitation du nombre de requêtes par intervalle de temps. Le chatbot
    intégré au site web est limité à 100 requêtes par jour par utilisateur pour
    maîtriser les coûts de l'API Claude.
\end{itemize}

Ces mesures de sécurité s'inscrivent dans une stratégie de défense en
profondeur qui couvre chaque couche de l'infrastructure~: conteneurisation
isolée, scan des vulnérabilités, gestion centralisée des secrets, et
protection des API exposées.
